\section{2016级工科数学分析下 B卷}

\subsection{资料来源}
\begin{itemize}
  \item 2016级工科数学分析试卷B.doc（试题骨架，DOC 文本抽取中大量公式对象缺失，仅作题型与题号参考）。
  \item 2016级工科数学分析期末试卷B解答.pdf（试题与答案，已按页面图和文本层人工校对）。
\end{itemize}

\subsection{试题}

\paragraph*{一、填空题（共 5 题，每题 2 分，共 10 分）}
\begin{enumerate}
  \item 由方程
  \[
    xyz+\sqrt{x^2+y^2+z^2}=\sqrt2
  \]
  所确定的函数 $z=z(x,y)$ 在点 $(1,0,-1)$ 处的全微分 $dz=$\underline{\hspace{4cm}}。

  \item 设 $\Gamma$ 为椭圆
  \[
    \frac{x^2}{9}+\frac{y^2}{4}=1
  \]
  取逆时针方向，则第二类曲线积分
  \[
    \oint_\Gamma x\,dy-y\,dx=
  \]
  \underline{\hspace{4cm}}。

  \item 向量场
  \[
    x^2yz\,\mathbf i+xy^2z\,\mathbf j+xyz^2\,\mathbf k
  \]
  在点 $(1,1,1)$ 的散度为\underline{\hspace{4cm}}。

  \item 初值问题
  \[
    \begin{cases}
      y'+\dfrac{3}{x}y=\dfrac{2}{x^3},\\
      y|_{x=1}=1
    \end{cases}
  \]
  的解为\underline{\hspace{4cm}}。

  \item 设 $f(x)=\pi x+x^2,\ -\pi<x<\pi$ 的傅里叶级数为
  \[
    \frac{a_0}{2}+\sum_{n=1}^{\infty}(a_n\cos nx+b_n\sin nx),
  \]
  则其中系数 $b_2=$\underline{\hspace{4cm}}。
\end{enumerate}

\paragraph*{二、单选题（共 5 题，每题 2 分，只有一个正确选项，共 10 分）}
\begin{enumerate}
  \item 二元函数 $f(x,y)$ 在点 $(a,b)$ 处的两个偏导数
  \[
    \frac{\partial f}{\partial x}(a,b),\qquad
    \frac{\partial f}{\partial y}(a,b)
  \]
  存在，则（\quad）。
  \begin{enumerate}
    \item[A.] $f(x,y)$ 在点 $(a,b)$ 处连续；
    \item[B.] $f(x,y)$ 在点 $(a,b)$ 处可微；
    \item[C.] $\displaystyle \lim_{x\to a}f(x,b)$ 和 $\displaystyle \lim_{y\to b}f(a,y)$ 都存在；
    \item[D.] $\displaystyle \lim_{(x,y)\to(a,b)}f(x,y)$ 存在。
  \end{enumerate}

  \item 已知曲面 $z=4-x^2-y^2$ 上点 $P$ 处的切平面平行于平面 $2x+2y+z=1$，则 $P$ 点的坐标是（\quad）。
  \begin{enumerate}
    \item[A.] $(1,-1,2)$；
    \item[B.] $(-1,1,2)$；
    \item[C.] $(1,1,2)$；
    \item[D.] $(-1,-1,2)$。
  \end{enumerate}

  \item 一个均匀物体由曲面 $z=x^2+y^2$ 及 $z=1$ 围成，则该物体的质心坐标为（\quad）。
  \begin{enumerate}
    \item[A.] $\left(\dfrac23,\dfrac23,\dfrac23\right)$；
    \item[B.] $\left(0,0,\dfrac23\right)$；
    \item[C.] $(0,0,1)$；
    \item[D.] $\left(\dfrac23,0,0\right)$。
  \end{enumerate}

  \item 微分方程
  \[
    (y-\ln x)\,dx+x\,dy=0
  \]
  是（\quad）。
  \begin{enumerate}
    \item[A.] 可分离变量方程；
    \item[B.] 一阶非齐次线性方程；
    \item[C.] 一阶齐次线性方程；
    \item[D.] 非线性方程。
  \end{enumerate}

  \item 下列级数条件收敛的是（\quad）。
  \begin{enumerate}
    \item[A.] $\displaystyle \sum_{n=1}^{\infty}\frac{(-1)^n}{n^2}$；
    \item[B.] $\displaystyle \sum_{n=1}^{\infty}\left(1-\cos\frac1n\right)$；
    \item[C.] $\displaystyle \sum_{n=1}^{\infty}\frac{n!}{n^n}$；
    \item[D.] $\displaystyle \sum_{n=1}^{\infty}\frac{(-1)^n}{n}$。
  \end{enumerate}
\end{enumerate}

\paragraph*{三、计算题（共 4 题，每题 7 分，共 28 分）}
\begin{enumerate}
  \item 设函数
  \[
    u(x,y)=y f\left(\frac{x}{y}\right)+x g\left(\frac{y}{x}\right),
  \]
  其中 $f$ 和 $g$ 具有连续的二阶导数，计算
  \[
    x\frac{\partial^2u}{\partial x^2}+y\frac{\partial^2u}{\partial x\partial y}.
  \]

  \item 计算二重积分
  \[
    \iint_D \frac{1-x^2-y^2}{1+x^2+y^2}\,dx\,dy,
  \]
  其中 $D$ 是区域 $x^2+y^2\le1,\ x\ge0,\ y\ge0$。

  \item 计算三重积分
  \[
    \iiint_\Omega z^2\,dx\,dy\,dz,
  \]
  其中 $\Omega$ 是平面 $x+y+z=1$ 与三个坐标面围成的区域。

  \item 计算第一类曲面积分
  \[
    \iint_\Sigma (x+y+z)\,dS,
  \]
  其中 $\Sigma$ 为平面 $y+z=5$ 被柱面 $x^2+y^2=25$ 所截得的部分。
\end{enumerate}

\paragraph*{四、解答题（共 4 题，每题 8 分，共 32 分）}
\begin{enumerate}
  \item 求方程
  \[
    y''+4y'+3y=0
  \]
  的一个解 $y=y(x)$，使其在点 $(0,2)$ 处与直线 $x-y+2=0$ 相切。

  \item 求圆柱面 $x^2+y^2=4,\ x\ge0,\ y\ge0$ 介于 $xOy$ 平面和曲面 $z=xy$ 之间的部分的面积。

  \item 设 $\Sigma$ 是锥面 $z=\sqrt{x^2+y^2}$ 被平面 $z=0$ 及 $z=1$ 截下的部分的下侧，计算第二类曲面积分
  \[
    \iint_\Sigma x\,dy\,dz+y\,dz\,dx+(z^2-2z)\,dx\,dy.
  \]

  \item 求幂级数
  \[
    \sum_{n=1}^{\infty}\frac{x^n}{n}
  \]
  的收敛域，并在收敛域上求其和函数。
\end{enumerate}

\paragraph*{五、证明题（共 2 题，每题 6 分，共 12 分）}
\begin{enumerate}
  \item 证明：曲面
  \[
    z=xe^{y/x}
  \]
  上所有点处的切平面都过一定点。

  \item 证明：函数项级数
  \[
    \sum_{n=0}^{\infty}x^n
  \]
  在 $(0,1)$ 上点态收敛，但不一致收敛。
\end{enumerate}

\paragraph*{六、应用题（共 1 题，共 8 分）}
从斜边长为 $l$ 的直角三角形中，求周长最大的直角三角形。

\subsection{答案与解析}

\paragraph*{一、填空题}
\begin{enumerate}
  \item 设
  \[
    F(x,y,z)=xyz+\sqrt{x^2+y^2+z^2}-\sqrt2.
  \]
  在点 $(1,0,-1)$ 处，
  \[
    F_x=\frac1{\sqrt2},\qquad F_y=-1,\qquad F_z=-\frac1{\sqrt2}.
  \]
  因此
  \[
    z_x=-\frac{F_x}{F_z}=1,\qquad z_y=-\frac{F_y}{F_z}=-\sqrt2,
  \]
  所以 $dz=dx-\sqrt2\,dy$。

  \item 由格林公式，
  \[
    \oint_\Gamma x\,dy-y\,dx=2S,
  \]
  其中椭圆面积 $S=\pi\cdot3\cdot2=6\pi$，故结果为 $12\pi$。

  \item
  \[
    \operatorname{div}(x^2yz,xy^2z,xyz^2)
    =2xyz+2xyz+2xyz=6xyz.
  \]
  在 $(1,1,1)$ 处取值为 $6$。

  \item 原方程乘以积分因子 $x^3$，得
  \[
    (x^3y)'=2.
  \]
  故 $x^3y=2x+C$。由 $y(1)=1$ 得 $C=-1$，所以
  \[
    y=\frac{2}{x^2}-\frac{1}{x^3}.
  \]

  \item
  \[
    b_n=\frac1\pi\int_{-\pi}^{\pi}(\pi x+x^2)\sin nx\,dx.
  \]
  由于 $x^2\sin nx$ 为奇函数，其积分为 $0$，故
  \[
    b_n=2\int_0^\pi x\sin nx\,dx=-\frac{2\pi(-1)^n}{n}.
  \]
  取 $n=2$，得 $b_2=-\pi$。
\end{enumerate}

\paragraph*{二、单选题}
\begin{enumerate}
  \item 两个偏导数存在只能保证沿坐标轴方向的一元极限存在，不保证二元连续或可微。选 C。
  \item 曲面法向量可取 $(2x,2y,1)$，与平面 $2x+2y+z=1$ 的法向量 $(2,2,1)$ 平行，得 $x=1,\ y=1,\ z=2$。选 C。
  \item 区域关于 $z$ 轴旋转对称，质心的 $x,y$ 坐标为 $0$；用柱坐标算得 $\bar z=\dfrac23$。选 B。
  \item 方程可化为 \(y'+\dfrac1x y=\dfrac{\ln x}{x}\)，是一阶非齐次线性方程。选 B。
  \item A 绝对收敛，B、C 收敛，D 是交错调和级数，只条件收敛。选 D。
\end{enumerate}

\paragraph*{三、计算题}
\begin{enumerate}
  \item 先求一阶偏导：
  \[
    u_x=f'\left(\frac{x}{y}\right)-\frac{y}{x}g'\left(\frac{y}{x}\right),
  \]
  \[
    u_y=-\frac{x}{y}f'\left(\frac{x}{y}\right)+g'\left(\frac{y}{x}\right).
  \]
  继续求导得
  \[
    u_{xx}=\frac1y f''\left(\frac{x}{y}\right)
    +\frac{y^2}{x^3}g''\left(\frac{y}{x}\right),
  \]
  \[
    u_{xy}=-\frac{x}{y^2}f''\left(\frac{x}{y}\right)
    -\frac{y}{x^2}g''\left(\frac{y}{x}\right).
  \]
  因此
  \[
    x u_{xx}+y u_{xy}=0.
  \]

  \item 作极坐标变换
  \[
    x=r\cos\theta,\qquad y=r\sin\theta,\qquad 0\le r\le1,\quad 0\le\theta\le\frac{\pi}{2}.
  \]
  则
  \[
    \iint_D \frac{1-x^2-y^2}{1+x^2+y^2}\,dx\,dy
    =
    \int_0^{\pi/2}d\theta\int_0^1\frac{1-r^2}{1+r^2}r\,dr.
  \]
  因为
  \[
    \int_0^1\frac{1-r^2}{1+r^2}r\,dr
    =\int_0^1\left(\frac{2r}{1+r^2}-r\right)\,dr
    =\ln2-\frac12,
  \]
  所以原积分为
  \[
    \frac{\pi}{4}(2\ln2-1).
  \]

  \item 对 $z$ 先积分外层，得
  \[
    \iiint_\Omega z^2\,dx\,dy\,dz
    =
    \int_0^1 z^2
    \left(\iint_{x+y\le1-z,\ x\ge0,\ y\ge0}dx\,dy\right)dz.
  \]
  截面面积为 $\dfrac12(1-z)^2$，因此
  \[
    \iiint_\Omega z^2\,dx\,dy\,dz
    =
    \frac12\int_0^1 z^2(1-z)^2\,dz
    =
    \frac1{60}.
  \]

  \item 曲面有显式方程 $z=5-y$，故
  \[
    dS=\sqrt{1+z_x^2+z_y^2}\,dx\,dy=\sqrt2\,dx\,dy.
  \]
  投影区域为 $x^2+y^2\le25$。于是
  \[
    \iint_\Sigma (x+y+z)\,dS
    =
    \iint_{x^2+y^2\le25}(x+5)\sqrt2\,dx\,dy.
  \]
  其中 $\iint x\,dx\,dy=0$，圆盘面积为 $25\pi$，故结果为
  \[
    125\sqrt2\,\pi.
  \]
\end{enumerate}

\paragraph*{四、解答题}
\begin{enumerate}
  \item 特征方程为
  \[
    \lambda^2+4\lambda+3=0,
  \]
  根为 $\lambda_1=-3,\ \lambda_2=-1$，故通解为
  \[
    y=C_1e^{-3x}+C_2e^{-x}.
  \]
  直线 $x-y+2=0$ 的斜率为 $1$，因此
  \[
    y(0)=2,\qquad y'(0)=1.
  \]
  即
  \[
    C_1+C_2=2,\qquad -3C_1-C_2=1.
  \]
  解得
  \[
    C_1=-\frac32,\qquad C_2=\frac72.
  \]
  所求解为
  \[
    y=-\frac32e^{-3x}+\frac72e^{-x}.
  \]

  \item 该柱面片可以看成以四分之一圆弧
  \[
    \Gamma:\ x^2+y^2=4,\quad x\ge0,\quad y\ge0
  \]
  为底线、以高度 $z=xy$ 竖直生成的曲面，故面积为
  \[
    S=\int_\Gamma xy\,ds.
  \]
  令
  \[
    x=2\cos t,\qquad y=2\sin t,\qquad 0\le t\le\frac{\pi}{2},
  \]
  则 $ds=2\,dt$，所以
  \[
    S=\int_0^{\pi/2}4\cos t\sin t\cdot2\,dt=4.
  \]

  \item 令 $\Sigma_1$ 为平面 $z=1$ 被锥面截下的圆盘，并取上侧，则 $\Sigma$ 与 $\Sigma_1$ 围成区域
  \[
    \Omega=\{(x,y,z):x^2+y^2\le z^2,\ 0\le z\le1\}.
  \]
  对向量场 $(P,Q,R)=(x,y,z^2-2z)$ 用高斯公式，有
  \[
    \iint_\Sigma P\,dy\,dz+Q\,dz\,dx+R\,dx\,dy
    +\iint_{\Sigma_1} P\,dy\,dz+Q\,dz\,dx+R\,dx\,dy
    =
    \iiint_\Omega 2z\,dx\,dy\,dz.
  \]
  右端为
  \[
    \int_0^1 2z\cdot \pi z^2\,dz=\frac{\pi}{2}.
  \]
  在 $\Sigma_1$ 上 $z=1$ 且取上侧，故
  \[
    \iint_{\Sigma_1} P\,dy\,dz+Q\,dz\,dx+R\,dx\,dy
    =
    \iint_{x^2+y^2\le1}(-1)\,dx\,dy=-\pi.
  \]
  因此所求积分为
  \[
    \frac{\pi}{2}-(-\pi)=\frac{3\pi}{2}.
  \]

  \item 幂级数的收敛半径为 $1$。当 $x=1$ 时，级数 $\sum 1/n$ 发散；当 $x=-1$ 时，级数
  $\sum (-1)^n/n$ 收敛。因此收敛域为 $[-1,1)$。

  记
  \[
    f(x)=\sum_{n=1}^{\infty}\frac{x^n}{n},\qquad x\in[-1,1).
  \]
  当 $x\in(-1,1)$ 时，
  \[
    f'(x)=\sum_{n=1}^{\infty}x^{n-1}=\frac1{1-x}.
  \]
  由 $f(0)=0$ 得
  \[
    f(x)=\int_0^x\frac{dt}{1-t}=-\ln(1-x),\qquad -1<x<1.
  \]
  在端点 $x=-1$ 由连续性延拓，故
  \[
    f(x)=-\ln(1-x),\qquad x\in[-1,1).
  \]
\end{enumerate}

\paragraph*{五、证明题}
\begin{enumerate}
  \item 设 $(x_0,y_0,z_0)$ 为曲面上一点，则
  \[
    z_0=x_0e^{y_0/x_0}.
  \]
  曲面的一个法向量为
  \[
    \left(-z_x,-z_y,1\right)\bigg|_{(x_0,y_0)}
    =
    \left(
      -\left(1-\frac{y_0}{x_0}\right)e^{y_0/x_0},
      -e^{y_0/x_0},
      1
    \right).
  \]
  于是切平面方程为
  \[
    -\left(1-\frac{y_0}{x_0}\right)e^{y_0/x_0}(x-x_0)
    -e^{y_0/x_0}(y-y_0)
    +(z-z_0)=0.
  \]
  将点 $(0,0,0)$ 代入左端，得到
  \[
    \left(1-\frac{y_0}{x_0}\right)x_0e^{y_0/x_0}
    +y_0e^{y_0/x_0}-z_0
    =
    x_0e^{y_0/x_0}-z_0=0.
  \]
  因此所有切平面都过定点 $(0,0,0)$。

  \item 对任意 $x\in(0,1)$，几何级数
  \[
    \sum_{n=0}^{\infty}x^n
  \]
  收敛，且和函数为
  \[
    S(x)=\frac1{1-x}.
  \]
  其部分和
  \[
    S_n(x)=\sum_{k=0}^{n}x^k=\frac{1-x^{n+1}}{1-x}.
  \]
  取 $x_n=1-\dfrac1n\in(0,1)$，则
  \[
    |S_n(x_n)-S(x_n)|
    =
    \frac{x_n^{n+1}}{1-x_n}
    =
    n\left(1-\frac1n\right)^{n+1}\not\to0.
  \]
  因此余项不能在 $(0,1)$ 上一致趋于 $0$，原函数项级数在 $(0,1)$ 上不一致收敛。
\end{enumerate}

\paragraph*{六、应用题}
设直角三角形两条直角边长分别为 $x,y$，则
\[
  x^2+y^2=l^2,\qquad x>0,\quad y>0,
\]
周长为
\[
  P=x+y+l.
\]
作拉格朗日函数
\[
  L(x,y,\lambda)=x+y+l-\lambda(x^2+y^2-l^2).
\]
驻点满足
\[
  1-2\lambda x=0,\qquad 1-2\lambda y=0,\qquad x^2+y^2=l^2.
\]
由前两式得 $x=y$，再由约束得
\[
  x=y=\frac{l}{\sqrt2}.
\]
此时周长最大，最大周长为
\[
  P_{\max}=(1+\sqrt2)l.
\]
